\documentclass[11pt]{article}
\title{AI Agent Reliability and Evaluation v1}
\date{}
\author{SIGMA AI-assisted research synthesis}
\usepackage[T1]{fontenc}
\usepackage[utf8]{inputenc}
\usepackage[margin=1in]{geometry}
\usepackage{amsmath,amssymb,array,booktabs,longtable,pdflscape}
\usepackage{microtype}
\usepackage{titlesec}
\emergencystretch=8em
\raggedbottom
\linespread{1.07}
\usepackage{xcolor}
\definecolor{blue}{HTML}{24485E}
\titleformat{\section}{\Large\bfseries\sffamily\color{blue}}{\thesection}{0.8em}{}
\titleformat{\subsection}{\large\bfseries\sffamily}{\thesubsection}{0.7em}{}
\setlength{\parskip}{.6em}
\setlength{\parindent}{0pt}
\begin{document}
\maketitle
\begin{center}\fbox{\parbox{.92\linewidth}{\textbf{Needs review: TeX verification is incomplete.} Review the archived proofreading diagnostics before relying on the typesetting.}}\end{center}
\begin{center}\small\textit{AI-generated research synthesis \textperiodcentered{} SIGMA}\end{center}

\begin{abstract}
The investigation asks: Autonomous AI agents are increasingly expected to complete complex tasks involving planning, reasoning, external tools, persistent state, and interaction with changing environments. The saved research reports a provisional finding, not independently verified here: investigator\_output/reasoning\_steps/0/argument: that a client lacks an acknowledgment. The modeled effect is the actual durable outcome, not an outbox insertion or dispatch intention. Actions outside this boundary and downstream delayed harms are not covered. A consequential recorded limitation or unsuccessful approach is: come/solution\_review/strategy\_findings/0/failure\_reason: Project or strategy changed; session paused before a new call.

The analysis distinguishes recorded arguments from demonstrated results, identifies remaining gaps, and does not claim experiments or proof steps absent from the frozen record.
\end{abstract}
\section*{About this paper}
This paper was generated by SIGMA, an AI-assisted research system, from the project's recorded investigations, findings, and supporting evidence.
The analysis and conclusions reflect an automated synthesis of that research. Although SIGMA incorporates evidence tracking, critical review, and consistency checks, the resulting paper has not necessarily undergone independent verification or peer review.
Readers should evaluate its arguments, supporting evidence, and stated limitations before relying on its conclusions.
Research style: Distinguish algorithmic guarantees from tests, benchmarks and conjectures. Document counterexamples, complexity assumptions and reproducibility.
\section*{Evidence coverage limitation}
The evidence index omitted 32187 material saved research field fragment(s) across 36 source(s). These fragments were not reviewed for this paper. Interpretations that could depend on omitted material remain provisional; the complete frozen source snapshot and omission ledger are retained with the generation record.
\textit{Claims discussed: }C1, C2
\section{Introduction and problem statement}
The investigation asks: Autonomous AI agents are increasingly expected to complete complex tasks involving planning, reasoning, external tools, persistent state, and interaction with changing environments. Demonstrations and benchmark results show that agents can perform many such tasks, but provide limited evidence that they remain dependable across repeated executions, unfamiliar conditions, long-running workflows, and settings where mistakes have significant consequences. This project asks how autonomous and semi-autonomous agents can be evaluated and engineered to deliver reliable, correct, recoverable, and reproducible results across extended, multi-step tasks.

Agent reliability is a system-level property. Probabilistic decisions, dynamically constructed plans, incomplete observations, tool interfaces, memory, orchestration, and external state all affect w

\textit{Claims discussed: }C1
\section{Methods and evidence}
This synthesis organizes the project's frozen research records and their saved evidence references. It does not claim new experiments.

\section{Background and research context}
Agent reliability is investigated here as dependable achievement of specified outcomes, rather than as the ability to produce an apparently successful execution. The research question concerns autonomous and semi-autonomous systems that combine planning, reasoning, external tools, persistent state, and interaction with changing environments. Its central distinction is between demonstrating that an agent can perform a task and establishing how dependably it satisfies that task\textquoteright{}s requirements under a defined execution protocol.

This framing makes the relationship among outcome measurement, failure recovery, and engineering intervention the organizing problem of the investigation. [C1; [Project overview]; [Strategy 1: Develop a rigorous framework for measuring AI agent reliability]]

Three complementary lines of inquiry define the scope. Measurement asks how correctness, robustness, and dependable behavior can be distinguished from superficial completion; failure analysis asks how faults propagate through extended execution and which outcomes remain detectable or recoverable. The engineering inquiry considers mechanisms such as constrained tool interfaces, explicit state and provenance, bounded retries, independent verification, escalation, and safe stopping as candidates for comparison, not as preferred solutions.

Together, these questions require attention both to what a system accomplishes and to the conditions under which a proposed safeguard could improve its outcomes. [C1; [Strategy 1: Develop a rigorous framework for measuring AI agent reliability]; [Strategy 2: Understand failure propagation and recovery in autonomous AI age]; [Strategy 3: Test architectural mechanisms for dependable agent outcomes]]

A concrete starting point is reliability under a fixed task distribution and execution protocol. The supplied formulation fixes a task distribution $D$, model and prompt versions, tool interfaces, environment-generation policy, and execution budget; a sampled task receives a prespecified batch of $r$ runs, from which an index $J$ is selected uniformly. If $Y_J$ indicates satisfaction of independently specified end-state requirements, safety constraints, and the deadline, the target is $p=\mathbb{E}[Y_J]$.

This target concerns success per attempted execution under that batch protocol: timeouts and crashes remain attempts, and the quantity is not success conditional on automated acceptance or success after unlimited retries. [C1; [Research session 1, artifact 1]]

The distinction between actual rubric satisfaction and an evaluator\textquoteright{}s acceptance decision is therefore foundational. The measurement formulations treat actual correctness, automated acceptance, and reference certification as separate variables, rather than defining correctness through agreement between evaluators. Reference validity is itself an evidential requirement: agreement alone cannot establish the proposed allowances for false certification or missed true successes.

The background question is consequently not only how to estimate a success rate, but also what outcome evidence makes that rate interpretable as reliability. [C1; [Research session 4, artifact 1]; [Research session 5, artifact 3]]

A deliberately narrow external-action model makes this problem concrete without presuming a general account of agent failure. Consider a task requiring exactly one irreversible external effect, with effect count $N$ and correctness condition $N=1$. The proposed recovery setting distinguishes a request lost before commitment from a committed request whose acknowledgment is lost, while assuming that both leave the agent with identical accessible observations.

This example directs attention to the boundary between local execution state and external outcomes: a completion record is an assertion to assess, not a substitute for inspecting the actual effect. [C1; [Research session 1, artifact 2]; [Research session 4, artifact 2]]

Even within this restricted setting, safety, completion, and truthful reporting require different definitions. Preventing multiple effects does not by itself require an effect to occur by a deadline, while verified completion additionally requires the stipulated evidence and reporting behavior. The temporal formulation also separates correctness at an observation horizon from truth when a success report was emitted; a missing or indeterminate report can fail a completion requirement without constituting a false success assertion.

These distinctions establish the vocabulary needed to assess recovery without treating harm containment, task achievement, and reporting correctness as interchangeable. [C1; [Research session 6, artifact 7]; [Research session 7, artifact 8]]

The engineering question is correspondingly comparative rather than architectural in the abstract. The proposed comparisons fix tasks, model and prompt versions, tool access, environment snapshots, and compute, time, and monetary caps, while separating the systems\textquoteright{} external state and randomizing execution order where necessary. Verification and recovery consume the same constrained resources as ordinary execution, and nonexecution or budget exhaustion remains in the evaluation denominator.

These controls define what evidence of an intervention\textquoteright{}s benefit would need to distinguish from differences in resources, starting conditions, or selective retention of attempts. [C1; [Research session 5, artifact 1]; [Research session 7, artifact 2]]

The analytical scope must nevertheless remain distinct from the broader research ambition. The supplied material provides restricted models, conditional reasoning, and proposed tests, not an executed comparison establishing architectural gains or benchmark-specific reliability. Its external-action boundary concerns actual durable effects rather than dispatch intentions, and does not cover unmodeled downstream consequences; correctness for one such operation is not end-to-end correctness for a dependent workflow.

The resulting foundation shifts the central question from whether an agent appears capable to which explicitly scoped outcomes, service conditions, and comparative evidence would justify calling its execution dependable. [C1; [Research session 7, artifact 1]; [Research session 6, artifact 6]; [Research session 5, artifact 2]]

\textit{Claims discussed: }C1
\section{Recorded findings and observations}
The recorded findings contribute a precise distinction between actual task correctness, observable completion signals, and the service conditions that make recovery safe. They are principally analytical results and proposed test specifications: the supplied evidence does not contain an implemented service audit, executed fault-injection study, or controlled agent comparison establishing an engineering benefit.

The configuration observation registered as C2 records a previous reasoning-effort setting of \textquotedblleft{}high\textquotedblright{}; it is not a measurement of correctness or evidence that additional reasoning improved reliability [C2; [Project overview]; [Research session 1]; [Research session 7]].

The measurement argument defines reliability relative to a fixed execution protocol rather than as an unrestricted property of an agent. A task is sampled from a specified distribution $D$, a prespecified batch of $r$ runs is executed, and an index $J$ is selected uniformly from that batch. With $Y_J=1$ denoting satisfaction of the independently specified end-state requirements, safety constraints, and deadline, the target is $p=\mathbb{E}[Y_J]$. This target concerns an attempted execution under the batch protocol, not success conditional on automated acceptance or after unlimited retries; timeouts and crashes remain attempted executions [[Research session 1, artifact 1]].

The central measurement deduction is that an automated pass rate does not identify this target. If $Z_J$ denotes automated acceptance, \$q=\textbackslash{}Pr(Z\_J=1)\$, \$f=\textbackslash{}Pr(Z\_J=1,Y\_J=0)\$, and \$g=\textbackslash{}Pr(Z\_J=0,Y\_J=1)\$, partitioning the possible outcomes gives $p=q-f+g$. Consequently, identical pass rates can coexist with entirely different actual correctness: an oracle that always accepts is compatible with executions that always succeed or always fail. A subsequent formulation makes reference error explicit: with automated acceptance $O$, reference certification $R$, \$a=\textbackslash{}Pr(O=1,R=0)\$, and an assumed allowance \$\textbackslash{}Pr(O=1,R=1,Y=0)\textbackslash{}leq\textbackslash{}eta\$, the recorded derivation gives $p\geq\max(0,q-a-\eta)$ without assuming independence among the three labels.

The allowance is an outcome-validity premise, not something established by agreement between judges [[Research session 1, artifact 1]; [Research session 4, artifact 1]].

For system comparisons, the recorded design treats the sampled task, rather than each repeated run, as the independent sampling unit. Both systems execute their scheduled batches on the same task in isolated copies of the initial environment; dependence among runs, systems, and judges within a task cluster is unrestricted. With prespecified stratum weights $w_h$ and system-specific correctness probabilities $p_{Ah}$ and $p_{Bh}$, the target contrast is $\Delta=\sum_h w_h(p_{Ah}-p_{Bh})$, under independent, representative task clusters.

This design preserves the distinction between repeated executions of one task and additional independently sampled tasks, while requiring versions, prompts, tools, horizons, and resource caps to remain fixed for the comparison [[Research session 5, artifact 1]; [Research session 5, artifact 3]].

A separate finite-roster argument changes the source of statistical uncertainty. It fixes the scheduled execution slots, their actual outcomes, and the observation horizon before random audit selection, retaining every scheduled slot under a prespecified nonexecution rule. The recorded argument uses hypergeometric inversion for marginal count bounds and a union bound for simultaneous coverage; its probability statement concerns audit selection conditional on the frozen roster, so it does not require independent execution outcomes.

This is a different inferential target from population reliability: its conservative use of possible-outcome sets depends on those sets containing the actual outcomes, and it supplies no automatic transport from the audited roster to deployment [[Research session 6, artifact 1]; [Research session 6]].

The recovery analysis identifies an information boundary that local reasoning cannot cross under the stipulated observations. One history loses a request before commitment, leaving no external effect; another commits the effect but loses its acknowledgment. If both expose identical information at recovery, and the only terminal actions are to skip or issue one retry that commits an additional effect, a local policy must use the same retry probability $q$ in both histories. Its respective error probabilities are $1-q$ and $q$, giving a sharp worst-case error bound of at least $1/2$, attained by $q=1/2$.

This is a bound for a restricted decision model, not a deployment failure estimate; in particular, the recorded correction preserves the retry-commitment premise because issuing a retry does not itself establish that another effect commits [[Research session 1, artifact 2]; [Research session 1]; [Research session 6, artifact 2]].

The sufficient recovery mechanism changes that service boundary rather than extracting additional information from the same ambiguous local state. It preserves a stable operation key and immutable payload, and atomically couples creation of the actual durable effect with creation of its completion record. Starting from empty effect and completion state, the first absent-key transition creates both, whereas a repeated matching request returns the retained result without another effect; payload mismatches are rejected. Under serialization, durability, noncolliding identities, and retention throughout every possible retry or delayed delivery, the recorded invariant establishes at-most-once effects.

At least one successful processing transition is additionally necessary for an effect to occur, and delivery of appropriate confirmation is separately necessary for verified completion [[Research session 4, artifact 2]; [Research session 2]].

The transition specifications make this distinction auditable by separating actual effects, completion records, pending messages, and volatile processing state. Sending changes the pending-message pool without creating an effect; delaying a message leaves it eligible for later processing, whereas dropping it removes it. In the refined model, receiving destructively removes a request into volatile execution state, so a pre-commit crash loses that consumed delivery rather than implicitly arranging redelivery. The displayed manual schedules preserve correspondence between effects and completion records under the atomic contract, but establish neither fairness nor eventual processing for arbitrary executions.

Their substantive contribution is therefore a determinate request-lifecycle specification, including outstanding work that can survive a client timeout [[Research session 4, artifact 2]; [Research session 4, artifact 4]; [Research session 4]].

Later specifications distinguish safety throughout execution from completion at a deadline. For one stable operation identity, the recorded finite-retention formulation defines safety as $C\leq1$ at every execution prefix, including after the evaluation deadline $H$, while effect completion requires $C(H)=1$. Verified completion further requires authoritative evidence and a truthful report by $H$; the safety condition alone permits permanent omission.

The modeled boundary is the actual durable outcome, not an outbox entry or dispatch intention, and does not cover effects outside that boundary or downstream delayed harms [[Research session 6, artifact 7]; [Research session 6, artifact 6]].

The reporting arguments add a corresponding evidential distinction: a correct state at the evaluation horizon does not by itself establish that a completion report was truthful when emitted. Report-time truth depends on commits preceding emission, and equal timestamp labels do not establish that ordering without suitable ordering evidence. Likewise, conformance to a conservative evidence-mapping rule implies actual-outcome containment only when the evidence used is truthful, complete where claimed, and correctly scoped and ordered.

Taken together, the recorded findings shift the reliability argument from obtaining plausible completion signals to specifying three separately accountable objects: the external outcome, the service contract governing its creation, and the evidence supporting the report [[Research session 7, artifact 8]; [Research session 7, artifact 9]; [Research session 7]].

\textit{Claims discussed: }C2
\section{Unsuccessful approaches and competing evidence}
The unsuccessful approaches narrow the reliability argument by showing where retry reasoning, duplicate suppression, and outcome labeling fail to establish dependable completion. The recorded preflight error documents an unsuccessful attempt, not evidence that a particular reliability mechanism is ineffective; no complete research artifact accompanies that attempt [C3; [Research session 2]]. The substantive negative evidence instead consists of explicit counterexamples and corrections to conditional arguments, rather than measured architectural regressions.

The retry argument concerns two histories that leave the agent with the same pending local state: in $W_0$, the original request produces no effect, whereas in $W_1$, it commits once but its acknowledgment is lost. For a task requiring exactly one irreversible effect by deadline $H$, let $q$ denote the probability of issuing a retry under this indistinguishable information. In the restricted construction where an issued retry commits by $H$, external-state correctness is bounded by $q$ in $W_0$ and by $1-q$ in $W_1$, yielding worst-case failure of at least $1/2$.

This is a conditional bound for the specified ambiguous-commitment decision, not a deployment failure-rate estimate or a bound covering authoritative lookup, compensation, or other distinguishing observations [[Research session 6, artifact 3]; [Research session 2]].

The attempted extension of that bound to arbitrary retry loss fails. An always-retry policy has $q=1$, but if every retry is dropped, the already-committed history $W_1$ retains exactly one effect; its external-state correctness probability is therefore $1$, contradicting the proposed upper bound $1-q=0$. The correction retracts the relaxed-delivery inequality while retaining the original assured-commitment construction.

It supplies no replacement bound for arbitrary delivery and adaptive observations, so retry issuance must not be treated as equivalent to another committed effect [[Research session 6, artifact 2]; [Research session 6, artifact 3]].

A separate overstatement concerns the claim that splitting effect creation from completion-record creation necessarily defeats at-most-once safety. Under the specified serialized, durable, no-repair contracts, effect-first execution can commit an effect, crash before recording completion, and then create a duplicate on a same-key retry. Marker-first execution has a different failure: a crash after the marker but before the effect causes subsequent same-key requests to return cached success while the effect count remains $C=0$.

Restricted marker-first suppression can thus preserve at-most-once safety while losing delivery and invalidating the meaning of its success record; the two write orders do not refute the same property [[Research session 6, artifact 3]; [Research session 7, artifact 3]; [Research session 7, artifact 5]].

Regenerating the operation key after ambiguous recovery does not provide a general repair. Although a fresh key can bypass the marker-first omission in the particular intermediate-crash trace, it also bypasses a valid completion record after a fully committed operation whose acknowledgment was lost. Even an implementation that atomically couples each effect with its receipt can then produce one effect under each of two keys, giving $C=2$ for the same logical operation.

The counterexample distinguishes per-key duplicate suppression from per-operation correctness and does not refute stronger split-write designs with durable repair mechanisms excluded by the model [[Research session 6, artifact 3]].

The atomic, stable-key alternative also supports a narrower conclusion than successful completion. Its conditional invariant prevents duplication without requiring delivery, but eventual effect creation additionally requires continued attempts and a service schedule in which an attempt completes; truthful terminal reporting further requires receipt delivery and client progress. Even those eventual-progress premises impose no finite completion deadline: useful deliveries can occur after any chosen $H$.

Accordingly, neither deadline success nor its probability follows from the safety invariant, and the manually specified traces establish no general fairness or eventual-processing property [[Research session 6, artifact 3]; [Research session 4]].

The outcome rubric contains a further failure that is independent of duplicate suppression: correct external state at the evaluation horizon does not establish truthful completion when the agent reports success. The challenged rubric permits a SUCCESS report by tick $15$ and evaluates effect state at $H=20$, allowing commitment after the report but before the horizon. That history can satisfy horizon-based correctness despite a premature completion claim. The correction therefore separates report-time truth from horizon-time state; it does not, by itself, change the predictions for the five specified schedules [[Research session 7, artifact 5]].

The subsequent temporal specification still leaves recoverable labeling gaps. Complete report capture together with exactly one matching effect row is insufficient to certify success if the captured report is INDETERMINATE, late, or accompanied by multiple final reports. Conversely, evidence of one timely matching SUCCESS does not establish uniqueness when capture is incomplete or report closure is absent. Definite success labels therefore require the full reporting contract\textemdash{}content, well-formedness, uniqueness, timeliness, and capture or closure\textemdash{}not merely the presence of a correct effect and a favorable report [[Research session 7, artifact 8]].

The proposed ordering ablation also fails to isolate uncertainty unless the retained evidence actually permits both event orders. If trusted timestamps on a common clock still place report emission at tick $15$ and commitment at tick $18$, removing a separate ordering certificate does not make their order unknown. A valid analytical contrast must specify which timestamp, clock-comparability, or ordering evidence is removed and show that both witness histories satisfy what remains. Otherwise, the ablation cannot support its claimed change in the outcome label [[Research session 7, artifact 8]].

These counterexamples reject particular inferences, not entire architectural families: issuing a retry need not create an effect, avoiding duplicates need not complete an operation, and obtaining a correct state by a deadline need not vindicate an earlier success report. The supplied record contains no inspected service or observer implementation, executed fault-injection comparison, or measured architectural benefit, so these corrections cannot establish net engineering improvement [[Research session 7, artifact 8]].

What they change is the burden of the central reliability argument: any claimed guarantee must identify the property preserved, the delivery and storage assumptions that sustain it, and the reporting evidence needed to distinguish genuine completion from a merely favorable final state.

\textit{Claims discussed: }C3
\section{Evidence limitations and uncertainties}
The evidence supports a distinction between conditional analytical results and demonstrated agent reliability, rather than treating the former as evidence of the latter. The central limitation is evidential: reported agreement does not independently verify an argument, its premises, or its applicability to an implemented system. The recorded analyses identify consequences of specified information, service, and evaluation contracts, but they do not establish that real agents and tools satisfy those contracts. This distinction governs how both positive arguments and negative findings can contribute to the paper\textquoteright{}s reliability claim [C4; [Research session 1]].

For recovery arguments, the consequential gap lies between the modeled information boundary and the actual external effect. The indistinguishability construction assumes that an uncommitted request and a committed request with a lost acknowledgment leave identical information available to the recovering policy. Timing, tool metadata, restored memory, or outstanding requests could defeat that equivalence in an implementation. Conversely, an apparently successful repair would not establish external correctness unless the effect oracle were independent and complete; no supplied implementation evidence establishes either the observation equivalence or these oracle properties [C4; [Research session 1]].

The sufficient service contracts have a similarly important boundary: the irreversible business effect must share the stipulated atomic, durable transition with its deduplication or authorization state. A key table, retained receipt, or atomic local intent is not by itself evidence that this boundary includes the actual downstream effect. Service atomicity, durability, retention, concurrency behavior, and preservation of operation identity by real agents remain unaudited.

The analytical conclusions therefore identify properties that would support restricted recovery guarantees, not measured compliance by existing services or a demonstrated causal improvement in end-to-end reliability [C4; [Research session 1]; [Research session 4]; [Research session 6]].

Evaluation uncertainty also concerns whether certification tracks truth, not merely whether a statistical calculation is correct. The population-comparison analysis assumes directional allowances for falsely accepted intervention outcomes and falsely rejected baseline successes, but those allowances lack independent outcome calibration. Likewise, conservative outcome sets contain the actual result only if the underlying snapshot, cutoff ordering, report capture, and nonexecution evidence are sound. Separate retention paths and archive digests do not establish those premises, and a silently false certificate can defeat containment despite correct application of the mapping rule.

Statistical error control cannot substitute for these missing connections between recorded evidence and actual outcomes [C4; [Research session 5]; [Research session 7]].

The statistical arguments nevertheless have different dependence requirements and should not be dismissed under a single independence objection. The population analysis assumes independent, representative task clusters while allowing dependence within a cluster; the finite-panel audit instead conditions on a frozen roster and locates randomness in uniform audit selection. Its coverage argument therefore does not require independent execution outcomes.

This distinction preserves the narrower finite-panel interpretation without establishing transport to a deployment population, and the recorded hand calculations and proposed coverage checks are not executed empirical audits or implementation tests [C4; [Research session 5]; [Research session 6]; [Research session 7]].

Negative findings must likewise be attached to the proposition they challenge. Although one recorded outcome is labeled as a reported counterexample, its detailed assessment treats the finite-panel lemma as convincing under its assumptions and locates the delivery objection in a different recovery claim. That objection restricts removal of the retry-commitment premise; it does not refute the audit lemma.

Historical references without substantive source content, incomplete artifacts, and unresolved prior obligations limit stronger compatibility claims, while an unavailable result supplies neither a successful demonstration nor a substantive refutation [C4; [Research session 6]; [Research session 5]; [Research session 3]; [Research session 2]].

Finally, the outcome being guaranteed is narrower than dependable task completion. At-most-once effect safety permits permanent omission, eventual successful delivery does not imply completion by a fixed deadline, and per-operation deduplication does not establish payload correctness, multi-service atomicity, or absence of delayed harm. No controlled evidence measures the associated recovery, completion, latency, storage, or monetary-cost tradeoffs.

The evidence thus changes the central argument from an assertion that a reliability mechanism works in practice to a more precise account of which contracts make a restricted guarantee possible and which outcome judgments remain unsupported [C4; [Research session 5]; [Research session 6]].

\begin{landscape}
\small
\begin{longtable}{@{}>{\raggedright\arraybackslash}p{\dimexpr(\linewidth-2\tabcolsep*3)/168*42\relax}>{\raggedright\arraybackslash}p{\dimexpr(\linewidth-2\tabcolsep*3)/168*42\relax}>{\raggedright\arraybackslash}p{\dimexpr(\linewidth-2\tabcolsep*3)/168*42\relax}>{\raggedright\arraybackslash}p{\dimexpr(\linewidth-2\tabcolsep*3)/168*42\relax}@{}}
\caption{Distinct inferential scopes in the recorded analyses. The limitations concern different bridges from a conditional argument to an applied reliability claim; they are not interchangeable objections.}\\
\toprule
\textbf{Analytical object} & \textbf{Conditional interpretation} & \textbf{Unestablished applicability requirement} & \textbf{Sources}\\
\midrule
\endfirsthead
\toprule
\textbf{Analytical object} & \textbf{Conditional interpretation} & \textbf{Unestablished applicability requirement} & \textbf{Sources}\\
\midrule
\endhead
\midrule
\multicolumn{4}{r}{\footnotesize Continued on next page}\\
\endfoot
\bottomrule
\endlastfoot
Ambiguous-commit recovery and stable-key repair & A restricted local-recovery error bound and sufficient per-operation effect contract. & Equivalent policy-visible observations; independent, complete effect observation; atomic durable coupling of the actual effect and retained record; preserved operation identity. & [Research session 1]\\
Fixed-protocol population comparison & A lower confidence bound for the success difference under fixed design, independent representative clusters, and directional truth-error allowances. & Outcome-grounded calibration, population matching, complete scheduled outcomes, isolation, and matched budgets in an actual evaluation. & [Research session 5]\\
Frozen-roster finite-panel audit & Coverage for a fixed roster under uniform independent audit selection and sound outcome sets, without execution-outcome independence. & Operationally compliant sampling and sound evidence mapping; neither deployment generality nor engineering benefit follows from the panel bound. & [Research session 6]; [Research session 7]\\
\end{longtable}
\par\smallskip{\footnotesize Evidence: [Research session 1], [Research session 5], [Research session 6], [Research session 7]; claims: C4}
\end{landscape}
\textit{Claims discussed: }C4
\section{Open questions and next steps}
The immediate research priority is to establish whether reference adjudication tracks actual task outcomes, rather than merely providing a second completion label. This is the unresolved obligation in C5: a reliability guarantee must either validate reference soundness or explicitly accommodate reference errors. The available work supplies conditional measurement arguments and proposed validation tests, but not an independently audited implementation demonstrating that the reference evidence is sound.

The next steps should therefore connect success certification to durable external outcomes before using certification rates to recommend an agent architecture. [C5; [Research session 1, artifact 1]; [Research session 6, artifact 3]]

The reference-error allowance provides a precise target for that validation. Let $Y$ denote actual rubric satisfaction, $O$ automated acceptance, and $R$ reference certification, with \$p=\textbackslash{}Pr(Y=1)\$, \$q=\textbackslash{}Pr(O=1)\$, and \$a=\textbackslash{}Pr(O=1,R=0)\$. The recorded conditional construction assumes \$c=\textbackslash{}Pr(O=1,R=1,Y=0)\textbackslash{}leq\textbackslash{}eta\$ and obtains $p\geq\max\{0,q-a-\eta\}$. The substantive next step is to justify an upper bound on this target-population joint miscertification probability; it is not interchangeable with a conditional error rate among certified executions.

Until independent outcome evidence supports such an allowance, results indexed by $\eta$ should be presented as sensitivity analysis rather than as an unconditional reliability certificate. [C5; [Research session 4, artifact 1]]

Comparative evaluations require attention to both directions of adjudication error. Although false reference rejection makes the preceding one-system lower bound conservative, the proposed system comparisons separately require justified directional allowances, particularly for missed true baseline successes. Agreement among judges does not establish either allowance.

A calibration study should therefore compare judgments with independently obtained outcome evidence in the intended task population, rather than treating agreement between automated and reference judgments as validation of truth. [C5; [Research session 4, artifact 1]; [Research session 5, artifact 1]; [Research session 5, artifact 3]]

For the ledger-backed tasks considered here, a useful pilot would first validate the outcome observer. It should capture actual durable effects by logical operation identity, including effects created under regenerated retry keys, rather than relying on completion records or dispatch intentions. It should also validate report capture and the ordering of commits, report events, and evaluation cutoffs; correctness at the horizon and truth at report emission should remain distinct outcomes. Missing or contradictory observations must not silently become either success or evidence that no effect occurred.

The table distinguishes the resulting validation questions so that success on one does not substitute for evidence on another. [C5; [Research session 6, artifact 3]; [Research session 7, artifact 3]; [Research session 7, artifact 8]]

Fault injection should then test whether the observer remains informative at the boundaries where recovery is ambiguous. The proposed harness must demonstrate that pre-commit request loss and post-commit acknowledgment loss expose identical recovery information to the agent while remaining distinguishable through independently audited service state. Tests should distinguish dropped requests from delayed requests, account for outstanding work, and include crashes between acknowledgment receipt and durable client completion. Delivery schedules must remain explicit: issuing a retry does not itself establish that another effect commits, as the dropped-retry counterexample shows.

These requirements make the experiment a test of external outcomes rather than a comparison of plausible transcripts. [C5; [Research session 1, artifact 2]; [Research session 4, artifact 4]; [Research session 6, artifact 3]]

The statistical implementation needs a separate validation programme. Proposed checks include exhaustive enumeration of small audit rosters, census and all-unknown cases, confidence-level sensitivity, correlated outcomes, and outcome-dependent evidence loss. Repeated-panel tests should also exercise task heterogeneity, within-task dependence, and asymmetric judge errors, while retaining every scheduled execution under a prespecified nonexecution rule.

Agreement in these finite tests would support correctness for the checked implementation cases; it would not prove general coverage or establish that the outcome sets supplied to the calculation contain the truth. [C5; [Research session 1, artifact 1]; [Research session 6, artifact 1]; [Research session 7, artifact 1]]

Once the reference evidence and calculation are validated, controlled agent comparisons can address engineering benefit. The proposed design holds tasks, model and prompt versions, tool access, environment snapshots, and execution budgets fixed; isolates system state; and randomizes condition order where appropriate. Comparisons among fencing, read-back-only retry, restart, and escalation should retain nonexecution and budget exhaustion, and report omissions, duplicates, false completion reports, unresolved outcomes, latency, and inclusive cost separately.

Deployment fault frequencies and asymmetric harm costs require separate estimation: neither conditional safety traces nor an experimentally selected fault mixture supplies those quantities. [C5; [Research session 7, artifact 2]; [Research session 5, artifact 2]]

Reference validation also needs an explicit response to failures of its own trust assumptions. Cross-execution and cross-build incidents should test the scope of certificate withdrawal and independently justified conditions for restoring certification; locally marking one observation unknown does not establish that other certificates remain sound. Broader applications additionally require outcome evidence for dependent multi-service workflows and delayed harms, rather than extending a single-operation ledger rubric by assertion.

These obligations change the central reliability argument from accepting independently named judges to requiring independently justified outcome evidence, with any remaining adjudication uncertainty carried into the reported guarantee. [C5; [Research session 7, artifact 6]; [Research session 7, artifact 9]; [Research session 4, artifact 2]]

\begin{longtable}{@{}>{\raggedright\arraybackslash}p{\dimexpr(\linewidth-2\tabcolsep*2)/112*28\relax}>{\raggedright\arraybackslash}p{\dimexpr(\linewidth-2\tabcolsep*2)/112*42\relax}>{\raggedright\arraybackslash}p{\dimexpr(\linewidth-2\tabcolsep*2)/112*42\relax}@{}}
\caption{Proposed validation layers for C5. Each layer addresses a different obligation; none of these tests is reported as an executed validation.}\\
\toprule
\textbf{Validation target} & \textbf{Proposed test and evidence} & \textbf{Distinction required when interpreting results}\\
\midrule
\endfirsthead
\toprule
\textbf{Validation target} & \textbf{Proposed test and evidence} & \textbf{Distinction required when interpreting results}\\
\midrule
\endhead
\midrule
\multicolumn{3}{r}{\footnotesize Continued on next page}\\
\endfoot
\bottomrule
\endlastfoot
Evidence-to-outcome mapping & Seal routine outputs before disclosing held-out reference evidence; exercise missing, conflicting, and same-tick observations, and validate report/commit ordering. & A mapper can follow its specification despite unsound input evidence. Mapper conformance and actual outcome correctness must be reported separately.\\
Service-contract conformance & Inject crashes and controlled delivery schedules; independently inspect actual effects, durable records, request rejection, and restart behavior. & A conditional invariant does not demonstrate implementation conformance. The treatment invariant must not itself serve as the outcome oracle.\\
Reference-observer soundness & Audit complete durable-effect capture, report receipts, cutoff ordering, namespace isolation, independent extraction, and nonexecution evidence; test cross-execution and cross-build trust incidents. & Synthetic reference agreement does not establish real-service soundness. Shared-root failures require a justified affected scope and certificate-withdrawal policy.\\
\end{longtable}
\par\smallskip{\footnotesize Evidence: [Research session 7, artifact 8], [Research session 7, artifact 9], [Research session 7, artifact 2], [Research session 7, artifact 6]; claims: C5}
\textit{Claims discussed: }C5
\section{Algorithms, results, and limitations}
\subsection{C1 -- Unresolved}
title: AI Agent Reliability and Evaluation v1
This is a recorded research item, not an independently verified conclusion; the cited saved evidence requires assessment.
\textit{Evidence: }[Project overview]
\subsection{C2 -- Unresolved}
events/4/project\_change/previous/reasoning\_effort: high
This is a recorded research item, not an independently verified conclusion; the cited saved evidence requires assessment.
\textit{Evidence: }[Project overview]
\subsection{C3 -- Unresolved}
events/10/failure\_kind: preflight\_error
This is a recorded research item, not an independently verified conclusion; the cited saved evidence requires assessment.
\textit{Evidence: }[Research session 2]
\subsection{C4 -- Limitation}
outcome/caveats/0: Agent-reported evidence is unverified; agreement is not independent verification.
This is a recorded research item, not an independently verified conclusion; the cited saved evidence requires assessment.
\textit{Evidence: }[Research session 1]
\subsection{C5 -- Unresolved}
investigator\_output/open\_obligations/0: Test reference-adjudication validity. If reference labels can falsely certify success, derive an additional error allowance or partial-identification bound rather than treating the present guarantee as unconditional.
This is a recorded research item, not an independently verified conclusion; the cited saved evidence requires assessment.
\textit{Evidence: }[Research session 1, artifact 1]
\section{Limitations}
\begin{itemize}
\item The saved record may contain unverified interpretations and unresolved evidence conflicts; these require independent assessment.
\end{itemize}
\section{Open obligations}
\begin{itemize}
\item investigator\_output/open\_obligations/0: Test reference-adjudication validity. If reference labels can falsely certify success, derive an additional error allowance or partial-identification bound rather than treating the present guarantee as unconditional.
\item investigator\_output/open\_obligations/0: Implement and audit the fault-injection harness, especially the assertion that the two histories expose identical information to the agent while the evaluator can distinguish their external states.
\item investigator\_output/open\_obligations/0: Recover and inspect original artifact [Research session 1, artifact 1] before asserting compatibility with its estimand, formulas, or saved local counterexample. One prior finding is omitted and the supplied recovery finding is truncated; historical uncertainty remains.
\item investigator\_output/open\_obligations/0: Next test: implement the specified transitions and replay T0\textendash{}T12, checking E/D invariants after every event, exact equality of T1/T2 agent-visible inputs at recovery, and complete accounting for pending requests. Independently inspect the harness observation projection and effect oracle before using its results.
\item investigator\_output/open\_obligations/0: Inspect the original of artifact [Research session 1, artifact 2] and the omitted finding before reconciling their complete assumptions and obligations; the present extract does not establish their resolution.
\item investigator\_output/open\_obligations/0: When execution becomes available, implement and independently audit the refined transitions, replay T0\textendash{}T12 with per-event delivery accounting and E/D checks, and audit the observation projection and effect oracle. In particular test crashes before RECEIVE, after RECEIVE, after PROCESS, and after ACK\_ENQUEUE against the stated loss semantics.
\item investigator\_output/open\_obligations/0: Independently justify or calibrate both directional truth-error allowances, especially baseline false rejection. Judge agreement alone is insufficient.
\item investigator\_output/open\_obligations/0: Execute the paired fault-injection experiment and independently audit actual service state, including missing acknowledgments, crashes at every receipt/effect boundary, and stable versus regenerated keys.
\item investigator\_output/open\_obligations/0: Validate target-population directional judge-error allowances, especially missed true baseline successes; agreement between judges cannot establish them.
\item investigator\_output/open\_obligations/0: Validate the implementation by exhaustive small-roster audit enumeration, including dependence, masking, census, and confidence-level sensitivity controls; no runtime is currently available.
\item investigator\_output/open\_obligations/0: Execute the crash-boundary test and independently audit actual effects, including client receipt-to-checkpoint crashes. The ledger obligation b963e848421fa9e188ec5db4683e603e35f508613827a711b79b0de2cc52b510 remains unresolved.
\item investigator\_output/open\_obligations/0: Implement and independently audit the six-cell fault-injection test, including concurrent requests, crashes during durable writes, and client completion persistence; neither runtime is currently available.
\item investigator\_output/open\_obligations/0: Execute the expiry and crash-boundary tests with independent external-effect auditing. Ledger b963e848421fa9e188ec5db4683e603e35f508613827a711b79b0de2cc52b510 and the related execution obligations remain unresolved.
\item investigator\_output/open\_obligations/0: Implement and independently audit the expiry-boundary test and the earlier six-cell crash test, including queued copies, durable-write failures, client completion persistence, and concurrent implementations.
\item investigator\_output/open\_obligations/0: Validate a concrete implementation's actual-effect atomicity, durable acknowledgment ordering, authoritative reads, failover, absent-key behavior, payload binding, and universal worker gating. The supplied execution and service-audit ledger duties remain unresolved.
\item investigator\_output/open\_obligations/0: Execute and independently audit the two-worker boundary schedule and prior crash tests, including paused-worker resumption, both GC orders, failover, durable-write failures, and client completion persistence. This addresses the next test for ledger obligations 223d7593f3c7d361a06e7d431e9fecb7580bfe0983c9e10b7555dab486ed5a86 and 73c690f66989c8781155d552241a1d2fbf18719
\end{itemize}
\section{Frozen obligation reconciliation}
Unlisted obligations are not presumed resolved. The archived research snapshot retains complete source references and the full reconciliation.
\begin{itemize}
\item Implement and audit the fault-injection harness, especially the assertion that the two histories expose identical information to the agent while the evaluator can distinguish their external states. \textit{ (Not listed by author; not presumed resolved; 10 saved source reference(s))}
\item Measure whether real agent policies preserve operation identity and payload across retries, checkpoint restoration, replanning, and concurrent workers. \textit{ (Not listed by author; not presumed resolved; 10 saved source reference(s))}
\item Test atomicity, crash durability, key retention, collision prevention, and payload-mismatch handling in any service claimed to satisfy the sufficient-repair assumptions. \textit{ (Not listed by author; not presumed resolved; 10 saved source reference(s))}
\item Estimate deployment frequencies and asymmetric costs of missing versus duplicate effects before choosing a practical retry policy; the worst-case bound does not supply these quantities. \textit{ (Not listed by author; not presumed resolved; 6 saved source reference(s))}
\item Evaluate authoritative lookup, escalation, compensation, and safe termination as additional policies, distinguishing harm containment from successful and verified task completion. \textit{ (Not listed by author; not presumed resolved; 10 saved source reference(s))}
\item Extend beyond one external action to dependent workflows with stale reads and partial multi-service commits; the current proof does not establish end-to-end workflow reliability. \textit{ (Not listed by author; not presumed resolved; 6 saved source reference(s))}
\item Run repeated, task-blocked comparisons with independent outcome audits and uncertainty estimates. Scripted protocol agreement alone would not establish general agent reliability. \textit{ (Not listed by author; not presumed resolved; 3 saved source reference(s))}
\item For open-ended tasks without an authoritative effect ledger, develop and validate an outcome oracle rather than transferring this experiment's observability assumptions. \textit{ (Not listed by author; not presumed resolved; 8 saved source reference(s))}
\item Implement and audit the fault-injection harness, including equality of all agent-accessible information and separation of evaluator-only external-state evidence. \textit{ (Not listed by author; not presumed resolved; 4 saved source reference(s))}
\item Measure operation-key and payload preservation by real agent policies across retry, checkpoint restoration, replanning, and concurrent execution. \textit{ (Not listed by author; not presumed resolved; 4 saved source reference(s))}
\item Test service atomicity, crash durability, retention, collision prevention, and payload-mismatch behavior. \textit{ (Not listed by author; not presumed resolved; 4 saved source reference(s))}
\item Estimate ambiguous-commit frequencies and asymmetric missing-effect and duplicate-effect costs before recommending a deployment policy. \textit{ (Not listed by author; not presumed resolved; 9 saved source reference(s))}
\item Evaluate authoritative lookup, escalation, compensation, and safe termination while separating harm containment from successful, verified completion. \textit{ (Not listed by author; not presumed resolved; 4 saved source reference(s))}
\item Extend the analysis to dependent workflows, stale observations, and partial multi-service commits. \textit{ (Not listed by author; not presumed resolved; 4 saved source reference(s))}
\item Conduct repeated, task-blocked comparisons with independent outcome audits and appropriate uncertainty estimates. \textit{ (Not listed by author; not presumed resolved; 9 saved source reference(s))}
\item Develop and validate outcome oracles for tasks without an authoritative external-effect ledger. \textit{ (Not listed by author; not presumed resolved; 4 saved source reference(s))}
\item Material, recoverable evidence gap: identical agent observations and independently observable external effects are assumptions of the proof, not properties demonstrated for an implemented harness. Timing, tool metadata, restored memory, and outstanding requests could invalidate the intended experimental pairing. \textit{ (Not listed by author; not presumed resolved; 7 saved source reference(s))}
\item Material, recoverable evidence gap: the sufficient repair assumes atomic, durable coupling of the business effect and deduplication record. A service that merely maintains a key table has not thereby satisfied that assumption. \textit{ (Not listed by author; not presumed resolved; 2 saved source reference(s))}
\end{itemize}
688 further distinct obligations are retained in the archived reconciliation.
\section{Conclusions}
Conclusions are limited to the findings and uncertainties documented in the cited saved research, as examined in the sections above.

\textit{Claims discussed: }C5
\appendix
\section{Selected source-linked research excerpts}
The following are bounded, selected excerpts of original research records. They are not exhaustive; consult the full frozen snapshot archived with this paper. They do not constitute independent verification.
\subsection{Research excerpt 1}
\textit{Source and field: }[Research session 6, artifact 4]: skeptic\_output/reasoning\_assessments\par
[\{"basis\_artifact\_ids": ["[Research session 6, artifact 2]", "[Research session 5, artifact 2]", "[Research session 6, artifact 3]", "[Research session 6, artifact 1]"], "disposition": "accepted", "objection": "", "rationale": "The supplied cycle 1 argument explicitly requires records to cover every possible retry and supplies the non-atomic and regenerated-key counterexamples. The peer architecture extract explicitly excludes expiry, while the measurement peer concerns a frozen roster.

None establishes finite-retention safety, and the current proof does not transfer their unsupported claims.", "step\_number": 1\}, \{"basis\_artifact\_ids": [], "disposition": "accepted", "objection": "", "rationale": "These are explicit premises of the conditional claim. Atomic checking with the actual effect, per-key ordering of deletion, crash durability, and monotonic retention time make the later concurrency and expiry arguments well-defined. They are not evidence that any deployed interface satisfies this model.", "step\_number": 2\}, \{"basis\_artifact\_ids": [], "disposition": "accepted", "objection": "", "rationale": "The first commit at 1 gives expiry at 11.

Deletion at 11 and an outstanding request's transition at 12 permit a second commit under the specified absent-record rule. Stopping sends at 3 neither removes nor disables that request. Thus C=2 follows without violating any expiry-only premise.", "step\_number": 3\}, \{"basis\_artifact\_ids": [], "disposition": "accepted", "objection": "", "rationale": "A request admitted at 2.5 can remain queued until 12, so an ingress\ldots{}

\subsection{Research excerpt 2}
\textit{Source and field: }[Research session 6, artifact 5]: skeptic\_output/reasoning\_assessments\par
[\{"basis\_artifact\_ids": ["[Research session 6, artifact 3]", "[Research session 5, artifact 2]", "[Research session 6, artifact 2]", "[Research session 6, artifact 1]"], "disposition": "accepted", "objection": "", "rationale": "The supplied cycle-1 argument explicitly requires nonexpiring receipts, so removing that premise requires a new invariant. The peer extracts likewise provide no finite-retention proof. The current ledger is nonempty, and the earlier wording error is explicitly corrected without claiming historical resolution.", "step\_number": 1\}, \{"basis\_artifact\_ids": [], "disposition": "accepted", "objection": "", "rationale": "These are explicit premises of the conditional claim.

Logical expiration at t=c+T, no replay refresh, and atomic actual-effect/receipt commitment specify the boundary behavior needed for the trace and proofs; none is presented as a demonstrated service property.", "step\_number": 2\}, \{"basis\_artifact\_ids": [], "disposition": "accepted", "objection": "", "rationale": "The delayed-copy and restart premises explicitly permit processing long after sending while preserving committed state.

Requiring eligibility and deduplication checks at the effect-producing transition excludes stale authorization from the conditional model.", "step\_number": 3\}, \{"basis\_artifact\_ids": [], "disposition": "accepted", "objection": "", "rationale": "At t=12 the receipt committed at t=1 is invalid because its expiry is 11. The ordinary service therefore takes its absent-receipt execution branch and increments C from one to two. Stable keys, atomic\ldots{}

\subsection{Research excerpt 3}
\textit{Source and field: }[Research session 7, artifact 1]: skeptic\_output/reasoning\_assessments\par
[\{"basis\_artifact\_ids": ["[Research session 6, artifact 1]"], "disposition": "accepted", "objection": "", "rationale": "The supplied extract of [Research session 6, artifact 1] is explicitly incomplete. Using it to identify recorded limitations while independently deriving the current result avoids importing unavailable proof steps. No retry-commitment inference is used.", "step\_number": 1\}, \{"basis\_artifact\_ids": [], "disposition": "accepted", "objection": "", "rationale": "Accepted as a proposed protocol premise, not an observed evaluation. Four schedules repeated four times define 16 pairs, and retaining all 32 scheduled executions with prespecified nonexecution failure keeps the denominator fixed.

The explicitly descriptive contrast does not assert a causal effect.", "step\_number": 2\}, \{"basis\_artifact\_ids": [], "disposition": "accepted", "objection": "", "rationale": "Accepted under the explicitly listed sound-certificate and well-defined-cutoff premises. The conjunction Y=E AND R distinguishes actual external success from reporting success. The step correctly treats certificate soundness as assumed rather than established by separate retention.", "step\_number": 3\}, \{"basis\_artifact\_ids": [], "disposition": "accepted", "objection": "", "rationale": "If each fact set contains its actual binary fact, the image under AND contains actual Y. Taking a Cartesian product preserves paired containment without component independence.

Widening conflicting facts avoids unsound empty intersections; a trusted duplicate state and a success report correctly imply failure\ldots{}

\subsection{Research excerpt 4}
\textit{Source and field: }[Research session 4, artifact 1]: skeptic\_output/reasoning\_assessments\par
[\{"basis\_artifact\_ids": [], "disposition": "accepted", "objection": "", "rationale": "Uniform selection within a fixed r-run batch defines a coherent marginal execution population. The localized unconditional error probability c differs correctly from a conditional reference-error rate. The allowance is accepted only as an explicit premise of the conditional claim.", "step\_number": 1\}, \{"basis\_artifact\_ids": [], "disposition": "accepted", "objection": "", "rationale": "Partitioning accepted failures by R gives Pr(O=1,Y=0)$\leq$a+c$\leq$a+η.

Accepted true successes form a subset of all true successes, so subtraction from q and clipping at zero establish the bound without judge independence.", "step\_number": 2\}, \{"basis\_artifact\_ids": [], "disposition": "accepted", "objection": "", "rationale": "For any feasible observed joint distribution, setting truth to failure outside the jointly accepted cell and allocating min(s,η) failure mass inside it attains max(0,s-η). Universal success is also feasible.

Mixtures preserve observations and the inequality on c, establishing every intermediate reliability value.", "step\_number": 3\}, \{"basis\_artifact\_ids": [], "disposition": "accepted", "objection": "", "rationale": "The two worlds have identical observed judge outputs but different truth labels and reliabilities. This establishes nonidentification from agreement alone, not an empirical benchmark failure rate or a contradiction to a separately justified η\textless{}1.", "step\_number": 4\}, \{"basis\_artifact\_ids": [], "disposition": "accepted", "objection": "", "rationale": "Averaging the run-specific\ldots{}

\subsection{Research excerpt 5}
\textit{Source and field: }[Research session 5, artifact 3]: skeptic\_output/reasoning\_assessments\par
[\{"basis\_artifact\_ids": [], "disposition": "accepted", "objection": "", "rationale": "As a premise of the explicitly conditional claim, fixed strata and weights together with uniform selection among scheduled runs define a coherent task-weighted estimand. Isolated paired copies and unrestricted within-cluster dependence do not conflict with independence between complete task clusters.", "step\_number": 1\}, \{"basis\_artifact\_ids": ["[Research session 4, artifact 1]"], "disposition": "accepted", "objection": "", "rationale": "The listed assumptions explicitly include these unconditional directional error bounds.

Intervention false certification can inflate intervention reliability, while baseline missed successes can depress baseline certification; these are the relevant adverse directions for a lower improvement bound. They are premises, not established calibration results. The supplied prior extract discusses only the lower-bound-oriented allowance.", "step\_number": 2\}, \{"basis\_artifact\_ids": [], "disposition": "accepted", "objection": "", "rationale": "The binary identity is exact. Dropping intervention missed successes and baseline false certifications gives the stated lower and upper bounds respectively.

Subtraction and nonnegative weighted summation preserve the inequality without judge independence.", "step\_number": 3\}, \{"basis\_artifact\_ids": [], "disposition": "accepted", "objection": "", "rationale": "With S\_A=1, S\_B=0, and both truth labels equal to one, the certification lower bounds differ by one while actual improvement is zero. Baseline missed-success\ldots{}

\subsection{Research excerpt 6}
\textit{Source and field: }[Research session 6]: steps/8/output/reasoning\_assessments\par
[\{"basis\_artifact\_ids": ["[Research session 6, artifact 2]", "[Research session 5, artifact 2]", "[Research session 6, artifact 3]", "[Research session 6, artifact 1]"], "disposition": "accepted", "objection": "", "rationale": "The supplied cycle 1 argument explicitly requires records to cover every possible retry and supplies the non-atomic and regenerated-key counterexamples. The peer architecture extract explicitly excludes expiry, while the measurement peer concerns a frozen roster.

None establishes finite-retention safety, and the current proof does not transfer their unsupported claims.", "step\_number": 1\}, \{"basis\_artifact\_ids": [], "disposition": "accepted", "objection": "", "rationale": "These are explicit premises of the conditional claim. Atomic checking with the actual effect, per-key ordering of deletion, crash durability, and monotonic retention time make the later concurrency and expiry arguments well-defined. They are not evidence that any deployed interface satisfies this model.", "step\_number": 2\}, \{"basis\_artifact\_ids": [], "disposition": "accepted", "objection": "", "rationale": "The first commit at 1 gives expiry at 11.

Deletion at 11 and an outstanding request's transition at 12 permit a second commit under the specified absent-record rule. Stopping sends at 3 neither removes nor disables that request. Thus C=2 follows without violating any expiry-only premise.", "step\_number": 3\}, \{"basis\_artifact\_ids": [], "disposition": "accepted", "objection": "", "rationale": "A request admitted at 2.5 can remain queued until 12, so an ingress\ldots{}

\subsection{Research excerpt 7}
\textit{Source and field: }[Research session 6, artifact 3]: skeptic\_output/reasoning\_assessments\par
[\{"basis\_artifact\_ids": ["[Research session 5, artifact 2]"], "disposition": "challenged", "objection": "Nonmaterial factual correction: five current ledger entries are supplied, so 'empty current ledger' is inaccurate. No historical resolution should be inferred either way; no rewrite cycle is needed for this wording alone.", "rationale": "The source is explicitly an incomplete extract, and declining to import its unshown proof is appropriate. Omitted history cannot establish resolution.

The only challenged component is the description of the current ledger.", "step\_number": 1\}, \{"basis\_artifact\_ids": ["[Research session 5, artifact 2]"], "disposition": "accepted", "objection": "", "rationale": "Under identical pre-retry information and assured retry commitment, success in the zero-effect history requires retry, while success in the one-effect history requires no retry. The stated bounds follow.

With every retry dropped, q=1 but the second history retains exactly one effect, directly disproving the relaxed second inequality.", "step\_number": 2\}, \{"basis\_artifact\_ids": [], "disposition": "accepted", "objection": "", "rationale": "These are explicit premises of the conditional claim, not established properties of deployed tools. They specify the durable states, replay behavior, serialization, operation identity, and excluded repair mechanisms sufficiently to determine the following traces.", "step\_number": 3\}, \{"basis\_artifact\_ids": [], "disposition": "accepted", "objection": "", "rationale": "Each table cell follows by applying the defined replay rule to the crash\ldots{}

\subsection{Research excerpt 8}
\textit{Source and field: }[Research session 7]: outcome/solution\_review/strategy\_findings/1/reasoning\_steps\par
[\{"argument": "Let C(k,t) count historical committed primary effects attributed to logical operation k; it is monotone and is not reduced by refunds or corrective records. The present incident domain permits a nonzero or corrupted initial state, unlike the clean-state progress domain described in the supplied [Research session 7, artifact 4] extract. A trustworthy post-fence read C\textgreater{}1 therefore exits that earlier guarantee's premises; it does not refute a result conditional on clean initialization. The original [Research session 7, artifact 2] is pointer-only here and has not been inspected.

Its reported invariant is not a premise of the new containment proof.", "basis\_artifact\_ids": ["[Research session 7, artifact 4]"], "epistemic\_status": "deduction", "number": 1, "research\_artifact\_id": "[Research session 7, artifact 7]", "review": \{"basis\_artifact\_ids": [], "disposition": "accepted", "objection": "", "rationale": "Defining C as historical committed effects makes compensation non-subtractive. A trustworthy duplicate count indicates that at least one premise of a clean-state safety guarantee is inapplicable; it need not specifically identify initialization as the failed premise.

The supplied ledger supports this scope comparison, not verification of the unavailable earlier proof."\}, "title": "Scope and provenance"\}, \{"argument": "Introduce Close(k,incident\_id), a proposed administrative primitive, not an existing property inferred from Advance. Close durably changes the operation's gate to CLOSED at a linearization point q. Every covered primary effect\ldots{}

\subsection{Research excerpt 9}
\textit{Source and field: }[Research session 7, artifact 5]: skeptic\_output/reasoning\_assessments\par
[\{"basis\_artifact\_ids": ["[Research session 7, artifact 3]", "[Research session 6, artifact 1]"], "disposition": "accepted", "objection": "", "rationale": "The supplied cycle-1 extract contains both the valid invariant and the challenged split-commit assertion. The current step explicitly withdraws the overstatement and does not import the truncated historical audit theorem.", "step\_number": 1\}, \{"basis\_artifact\_ids": [], "disposition": "accepted", "objection": "", "rationale": "These are proposed premises, not inspected implementation facts.

The E schema permits duplicate request identifiers, D supplies logical-key uniqueness, and the text explicitly identifies linked-row consistency as an additional application obligation rather than falsely attributing it to the foreign key.", "step\_number": 2\}, \{"basis\_artifact\_ids": [], "disposition": "accepted", "objection": "", "rationale": "Under the conditional claim, atomic rollback includes E, ambiguous commits are resolved through a new D lookup, and restart never blindly resumes an old append.

These requirements close the specified duplicate-creation paths without asserting an actual database satisfies them.", "step\_number": 3\}, \{"basis\_artifact\_ids": ["[Research session 7, artifact 3]"], "disposition": "accepted", "objection": "", "rationale": "The first successful absent-key transaction creates both rows. D's unique key forces competing transactions to abort their E writes, and retained D prevents later insertion. Complete-state crash recovery preserves the invariant; a safety violation in an infinite execution\ldots{}

\subsection{Research excerpt 10}
\textit{Source and field: }[Research session 7, artifact 7]: skeptic\_output/reasoning\_assessments\par
[\{"basis\_artifact\_ids": [], "disposition": "accepted", "objection": "", "rationale": "Defining C as historical committed effects makes compensation non-subtractive. A trustworthy duplicate count indicates that at least one premise of a clean-state safety guarantee is inapplicable; it need not specifically identify initialization as the failed premise. The supplied ledger supports this scope comparison, not verification of the unavailable earlier proof.", "step\_number": 1\}, \{"basis\_artifact\_ids": [], "disposition": "accepted", "objection": "", "rationale": "This is an explicit premise of the conditional claim.

Atomic ordering of closure with every covered effect creation, persistent CLOSED state, and prohibition of ordinary reopening are sufficient contract ingredients for the later invariant; no implementation is inferred.", "step\_number": 2\}, \{"basis\_artifact\_ids": [], "disposition": "accepted", "objection": "", "rationale": "These are additional stated authority, persistence, and evidence premises, not established capabilities. Durable fail-closed dispatch suppression supports restart behavior, while certificate identity and ordering requirements prevent confusing a requested close with a confirmed close.

Outbox persistence alone is correctly not treated as delivery.", "step\_number": 3\}, \{"basis\_artifact\_ids": [], "disposition": "accepted", "objection": "", "rationale": "The proposed incident action suppresses new autonomous business dispatches without claiming cancellation of already pending effects. Its distinction between immediate local stopping and suppression\ldots{}

\subsection{Research excerpt 11}
\textit{Source and field: }[Research session 7, artifact 9]: skeptic\_output/reasoning\_assessments\par
[\{"basis\_artifact\_ids": [], "disposition": "challenged", "objection": "[Research session 6, artifact 1] is supplied only as a pointer, not the substantive extract described in this step. Use the ledger as a recorded constraint, not as an inspected proof.", "rationale": "The continuity restriction and preservation of unknown duties are appropriate, but the claimed inspection exceeds the supplied evidence.", "step\_number": 1\}, \{"basis\_artifact\_ids": [], "disposition": "accepted", "objection": "", "rationale": "This is an explicit standalone rubric, conditional on the listed well-defined isolated history premise.

Exact-one commit and exact-one valid final receipt make additional events failures, and historical before-H labels are not rewritten by later events.", "step\_number": 2\}, \{"basis\_artifact\_ids": [], "disposition": "accepted", "objection": "", "rationale": "The specified branches preserve nonempty component sets. Completeness is required for exact-one positive conclusions, while correctly scoped pre-H irreversible violations can establish failure without complete history.

Conflict widening is conservative rather than an empty-set inference.", "step\_number": 3\}, \{"basis\_artifact\_ids": [], "disposition": "accepted", "objection": "", "rationale": "M compares output with the input-version specification; V compares certificate assertions with history; Z tests membership of the actual outcome; B depends on review-time trust information. These definitions do not equate ignorance with truth or a hold with task failure.", "step\_number": 4\}, \{"basis\_artifact\_ids": [],\ldots{}

\subsection{Research excerpt 12}
\textit{Source and field: }[Research session 1]: outcome/solution\_review/proof\_obligations\par
["Analytically discharged within the conditional model, cycle 1, current Investigator reasoning entries 1\textendash{}3: construct indistinguishable histories and establish the sharp minimax error bound.", "Analytically discharged within the conditional model, cycle 1, current Investigator reasoning entry 4: distinguish conditional average error from worst-case error and overall deployment failure.", "Analytically discharged within the conditional model, cycle 1, current Investigator reasoning entries 5\textendash{}6: explain why local rollback is insufficient and why atomic durable same-key deduplication preserves exactly one effect.", "Open: establish that a concrete harness implements the assumed information boundary and transition semantics.", "Open: establish that concrete agent policies and services satisfy the repair assumptions.", "Open: measure controlled agent-level outcomes and establish the applicability limits for deployment and multi-step workflows."]

\subsection{Research excerpt 13}
\textit{Source and field: }[Research session 1, artifact 1]: skeptic\_output/reasoning\_assessments\par
[\{"basis\_artifact\_ids": [], "disposition": "accepted", "objection": "", "rationale": "The uniform run selection makes E[Y\_J] the expected within-batch success fraction under the fixed protocol. Explicitly including crashes and timeouts avoids conditioning on successful execution or oracle acceptance. This defines an estimand rather than establishing deployment representativeness.", "step\_number": 1\}, \{"basis\_artifact\_ids": [], "disposition": "accepted", "objection": "", "rationale": "Partitioning the joint binary outcomes gives q=P(Y=1,Z=1)+f and p=P(Y=1,Z=1)+g, so p=q-f+g.

The two constant-oracle examples correctly demonstrate nonidentification from q alone, and f is correctly distinguished from a conditional false-positive rate.", "step\_number": 2\}, \{"basis\_artifact\_ids": [], "disposition": "accepted", "objection": "", "rationale": "Under the listed independent-cluster assumptions, each X\_i lies in [0,1] and has expectation q because uniform run selection averages the r oracle decisions. One-sided Hoeffding therefore has exponent 2Nt\textasciicircum{}2.

Dependence among runs within X\_i does not invalidate that application or justify an Nr exponent.", "step\_number": 3\}, \{"basis\_artifact\_ids": [], "disposition": "accepted", "objection": "", "rationale": "Distribution-matched independent audit clusters, uniform run selection, and correct adjudication make the A\_l independent Bernoulli(f). For k\textless{}m, the binomial lower-tail probability decreases continuously from 1 to 0 as u increases, giving the stated unique root. Inversion yields the conservative one-sided upper confidence limit, including\ldots{}

\subsection{Research excerpt 14}
\textit{Source and field: }[Research session 1, artifact 2]: objective\_review/proof\_obligations\par
["Analytically discharged within the conditional model, cycle 1, current Investigator reasoning entries 1\textendash{}3: construct indistinguishable histories and establish the sharp minimax error bound.", "Analytically discharged within the conditional model, cycle 1, current Investigator reasoning entry 4: distinguish conditional average error from worst-case error and overall deployment failure.", "Analytically discharged within the conditional model, cycle 1, current Investigator reasoning entries 5\textendash{}6: explain why local rollback is insufficient and why atomic durable same-key deduplication preserves exactly one effect.", "Open: establish that a concrete harness implements the assumed information boundary and transition semantics.", "Open: establish that concrete agent policies and services satisfy the repair assumptions.", "Open: measure controlled agent-level outcomes and establish the applicability limits for deployment and multi-step workflows."]

\subsection{Research excerpt 15}
\textit{Source and field: }[Research session 4]: outcome/solution\_review/proof\_obligations\par
["Lifecycle and request conservation: cycle 2, steps 1\textendash{}4 establish explicit location transfers and crash outcomes under the chosen transport premises.", "Durable atomic safety: cycle 2, step 5 establishes preservation under the refined transitions; cycle 1, step 3 in artifact [Research session 4, artifact 2] supplies the compatible saved inductive argument.", "Affected finite traces: cycle 2, steps 6\textendash{}10 supply consistent per-event T5, T8, and both T12 tables, including consumed deliveries and acknowledgments.", "Safety versus existence: cycle 1, step 4 in artifact [Research session 4, artifact 2] requires a completed matching PROCESS for exactly-one existence; cycle 2 adds no liveness guarantee.", "Naive split-write failure boundary: cycle 1, step 7 in artifact [Research session 4, artifact 2] retains valid local duplicate-effect and false-completion counterexamples; these do not refute all non-atomic recovery protocols.", "Model-to-system correspondence, real-agent identity preservation, controlled causal benefit, broader workflow applicability, and historical reconciliation remain unestablished."]

\subsection{Research excerpt 16}
\textit{Source and field: }[Research session 2]: outcome/conclusion\par
"\#\#\# Unresolved investigation\textbackslash{}n\textbackslash{}nThe session is queued and has not saved a complete Investigator/Skeptic finding. No answer to the project question is established yet.\textbackslash{}n\textbackslash{}n\#\#\# Findings across strategies\textbackslash{}n- **Develop a rigorous framework for measuring AI agent reliability \textemdash{} unresolved.**\textbackslash{}n- **Understand failure propagation and recovery in autonomous AI agents \textemdash{} unresolved.**\textbackslash{}n- **Test architectural mechanisms for dependable agent outcomes \textemdash{} unresolved.**\textbackslash{}n\textbackslash{}n\#\#\# Session caveats\textbackslash{}n- No complete research artifact is available.\textbackslash{}n\textbackslash{}n\#\#\# Verification and scope\textbackslash{}n\textbackslash{}nVerification: Incomplete.\textbackslash{}n\textbackslash{}nThis Conclusion summarizes agent-reported findings for this session.

SIGMA has not independently verified the result; inspect the detailed Outcome and saved evidence before relying on it."

4 further research records are retained in the archived source snapshot.
\section*{Evidence and authorship note}
Snapshot SHA-256: 080d3b70b67cfbe346acdb0f565d4a61c09ec4819229f2dd92dded43d9de40d9
Generated by an AI author from the cited saved research snapshot. This document is not independent verification. No external bibliography was added.
\section*{Validation warnings}
\begin{itemize}
\item Evidence indexing omitted 32187 material saved research field fragment(s). Those fragments were not reviewed for this paper; the interpretation is provisional where the omitted evidence may matter. Consult the frozen source snapshot and omission ledger for full details.
\item The AI planning stage did not yield a usable outline. A conservative, source-linked outline was reconstructed from saved evidence without repeating the planning request; review the resulting paper's structure.
\item The evidence index explicitly omitted 32191 saved field fragment(s); see the generation coverage ledger and review the complete archived snapshot.
\item Structured paper indexing selected 437 source-linked text fragments from 36 archived sources without AI extraction calls. The complete frozen snapshot remains archived; this bounded index does not constitute review of every saved field.
\item The frozen evidence contains open proof/research obligation(s) requiring reconciliation: 690 item(s) are not listed by the author and are not presumed resolved; 16 author-listed item(s) have no exact frozen-record match. See the obligation reconciliation.
\item The methods narrative has no linked claim IDs; check it for unsupported factual statements.
\item AI-authored analysis is not independent formal verification; model agreement does not establish proof.
\end{itemize}
\section*{Typesetting verification note}
Needs review: TeX verification is incomplete. The authored research content and safe accepted typesetting corrections were retained. Review the saved proofreading diagnostics before relying on the layout.
\end{document}
